\documentclass[11pt]{article}
\usepackage[T1]{fontenc}
\usepackage{lmodern}
\usepackage[margin=1in]{geometry}
\usepackage{amsmath,amssymb,amsthm}
\usepackage{microtype}
\usepackage[hidelinks]{hyperref}
\hypersetup{pdftitle={A two-factor criterion for unimodality of stretched q-integer products},
  pdfauthor={Xiaojian Zeng}}
\newtheorem{theorem}{Theorem}

\newcommand{\ind}{\mathbf{1}}
\title{A two-factor criterion for unimodality\protect\\
of stretched $q$-integer products}
\author{Xiaojian Zeng\thanks{PBC School of Finance, Tsinghua University,
100190. Email: \href{mailto:xiaojianzeng@foxmail.com}{xiaojianzeng@foxmail.com}.}}
\date{September 14, 2026}

\begin{document}
\maketitle
\begin{abstract}
For every integer $r\ge2$ and positive integers $a,c,b$, we prove that
$[a]_q[c]_q[b]_{q^r}$ is weakly unimodal exactly when $r$ divides $a$ or $c$,
or $b\le1+\lfloor a/r\rfloor+\lfloor c/r\rfloor$.
This establishes the two-factor case of Conjecture~5.4 of Connelly, Ito,
Martinez, Shevchenko, and Yang, extending their stretch-$2$ characterization
to arbitrary stretch. The proof compares two counts of multiples of $r$
and controls their endpoints by a uniform midpoint estimate. When the
criterion fails, an explicit first-half coefficient difference equals $-1$.
The original conjecture in its full parameter range remains unresolved here.
\end{abstract}

\section{Statement and context}
For a positive integer $m$, write $[m]_q=1+q+\cdots+q^{m-1}$, with $q$
an indeterminate. A polynomial $P(q)=\sum_{n=0}^D p_nq^n$ is
\emph{weakly unimodal} if, for some $m$,
$p_0\le\cdots\le p_m\ge\cdots\ge p_D$.
The sequence includes every integer degree, including degrees with zero
coefficient. We call $P$ symmetric if $p_n=p_{D-n}$.

In their study of $q$-Fibonomial coefficients, Connelly, Ito, Martinez,
Shevchenko, and Yang proposed a criterion for products with one stretched
factor \cite[Conjecture~5.4]{CIMSY2026}. For $r\ge2$, $k\ge1$, and positive
integers $a_1,\ldots,a_k,b$, their conjecture asserts that
\begin{equation}\label{eq:parent}
 r\mid a_i\text{ for some }i
 \quad\text{or}\quad
 b\le1+\sum_{i=1}^k\left\lfloor \frac{a_i}{r}\right\rfloor
\end{equation}
is sufficient for the unimodality of
$[a_1]_q\cdots[a_k]_q[b]_{q^r}$, and is also necessary when $k\le3$ or
$r\le3$. We establish precisely its $k=2$ restriction.

\begin{theorem}\label{thm:main}
Let $r\ge2$ be an integer and let $a,c,b$ be positive integers. Then
$[a]_q[c]_q[b]_{q^r}$ is weakly unimodal if and only if
\begin{equation}\label{eq:criterion}
 r\mid a\quad\text{or}\quad r\mid c\quad\text{or}\quad
 b\le1+\left\lfloor \frac{a}{r}\right\rfloor+\left\lfloor \frac{c}{r}\right\rfloor.
\end{equation}
The alternatives are inclusive, and parameters equal to $1$ are allowed.
\end{theorem}

The divisibility branch is already known
\cite[Lemma~4.2 and Corollary~4.3]{CIMSY2026}, as is the case of one ordinary
factor \cite[Proposition~5.3]{CIMSY2026}. The closest two-factor result is
\cite[Proposition~4.5]{CIMSY2026}: at $r=2$ the criterion is $2b\le a+c$
or at least one of $a,c$ is even. For two odd lengths, this is exactly
\eqref{eq:criterion}. Their proof compares intervals contributing $+1$
and $-1$ to a coefficient difference. We use the same reduction; the step
that permits every stretch is the floor estimate in
\eqref{eq:floor-bound}--\eqref{eq:residue}. The proof below includes the
known boundary cases to make the argument self-contained. It makes no
claim of priority beyond this comparison with the inspected sources.

\section{Proof of the criterion}
We first recall the classical closure of symmetric, nonnegative, unimodal
polynomials under multiplication \cite{Andrews1975}. Here is a direct
argument. If $F(q)=\sum_{i=0}^E f_iq^i$ has these properties, then
\begin{equation}\label{eq:centered}
 F(q)=\sum_{j=0}^{\lfloor E/2\rfloor}
 (f_j-f_{j-1})q^j[E-2j+1]_q,\qquad f_{-1}=0.
\end{equation}
This follows by telescoping up to the midpoint and then using symmetry.
All the displayed weights are nonnegative. The product $[u]_q[v]_q$ has
symmetric trapezoidal coefficients, since its coefficient at degree $n$
counts pairs $(i,j)$ with $i+j=n$, $0\le i<u$, and $0\le j<v$.
Expanding two decompositions \eqref{eq:centered} therefore expresses their
product as a nonnegative sum of shifted trapezoids with a common center.
Their sum is symmetric and unimodal.

\begin{proof}[Proof of Theorem~\ref{thm:main}]
If $a=rA$ for a positive integer $A$, then
\[
 [a]_q[b]_{q^r}=[r]_q\left.\bigl([A]_x[b]_x\bigr)\right|_{x=q^r}.
\]
Its coefficients repeat each coefficient of the symmetric trapezoid
$[A]_x[b]_x$ exactly $r$ times. Hence it is symmetric and unimodal, and
the closure just proved permits multiplication by $[c]_q$.
The case $r\mid c$ is identical.

Assume henceforth that neither $a$ nor $c$ is divisible by $r$.
By commutativity we may suppose $a\le c$. Put
\[
 P(q)=[a]_q[c]_q[b]_{q^r}=\sum_{n=0}^D p_nq^n,
 \qquad D=a+c-2+r(b-1),
\]
and set $p_n=0$ outside $[0,D]$. Each factor is symmetric, so
unimodality is equivalent to
\begin{equation}\label{eq:half}
 \Delta_n:=p_n-p_{n-1}\ge0
 \qquad\left(1\le n\le\left\lfloor \frac D2\right\rfloor\right).
\end{equation}
When $D$ is odd the two central coefficients agree by symmetry; when
$D=0$ the comparison range is empty.

For integers $L,U$, define
\[
 N(L,U)=\#\{j\in\mathbb Z:0\le j\le b-1,\ L\le rj\le U\}.
\]
Taking coefficients in
$(1-q)[a]_q[c]_q=[a]_q-q^c[a]_q$ gives the identity
\begin{equation}\label{eq:counts}
 \Delta_n=N(n-a+1,n)-N(n-a-c+1,n-c).
\end{equation}
In particular, $\Delta_n\ge0$ for $n<c$, because the second count vanishes.
Write
\begin{equation}\label{eq:parameters}
 \begin{gathered}
 a=rA+s,\qquad c=rC+t,\qquad A,C\ge0,\quad 1\le s,t\le r-1,\\
 M=A+C,\qquad B=b-1,\qquad D=r(M+B)+s+t-2.
 \end{gathered}
\end{equation}

\emph{Sufficiency.}
Suppose $B\le M$. Only $c\le n\le\lfloor D/2\rfloor$ remains to be
considered. If this range is nonempty then $B\ge1$, since for $B=0$ one
has $D/2=(a+c-2)/2\le c-1$.
Since $rB\le rM\le a+c-2$, the midpoint bound gives $n\le a+c-2$.
Moreover,
\[
 0\le n-c\le\frac{a-c-2+rB}{2}\le\frac{rB-2}{2}<rB,
 \qquad n-a\ge0.
\]
Thus the second count in \eqref{eq:counts} starts at $j=0$ and its upper
endpoint is below $B$; the first count has lower endpoint at least $1$.
Accounting for the possible upper clipping of the first count yields
\begin{equation}\label{eq:clipped}
 \Delta_n=\max\!\left\{0,\min\!\left(B,\left\lfloor \frac nr\right\rfloor\right)
                -\left\lfloor \frac{n-a}{r}\right\rfloor\right\}
               -\left\lfloor \frac{n-c}{r}\right\rfloor-1.
\end{equation}
Set
\[
 F=\left\lfloor \frac{n-a}{r}\right\rfloor+\left\lfloor \frac{n-c}{r}\right\rfloor+1,
 \qquad h=\left\lfloor \frac nr\right\rfloor.
\]
We will show $F\le\min(B,h)$. Floor subadditivity and $2n\le D$ imply
\begin{equation}\label{eq:floor-bound}
 F\le\left\lfloor \frac{2n-a-c}{r}\right\rfloor+1
   \le\left\lfloor \frac{rB-2}{r}\right\rfloor+1=B.
\end{equation}
For the other bound, write $n=rh+v$, $0\le v<r$, and put
$e=\ind_{v<s}+\ind_{v<t}$, where $\ind$ denotes an indicator. Then
\begin{equation}\label{eq:residue}
 F=2h-M+1-e.
\end{equation}
The midpoint estimate gives
\[
 n\le\frac D2\le rM+\frac{s+t-2}{2}<r(M+1),
\]
so $h\le M$. If $h\le M-1$, equation~\eqref{eq:residue} gives $F\le h$.
If $h=M$, the same estimate gives $2v\le s+t-2$, whence
$v<\max(s,t)$ and $e\ge1$; again $F\le h$.
It follows that
\[
 \min(B,h)-\left\lfloor \frac{n-a}{r}\right\rfloor
 \ge\left\lfloor \frac{n-c}{r}\right\rfloor+1\ge1.
\]
Equation~\eqref{eq:clipped} now yields $\Delta_n\ge0$, proving sufficiency
by symmetry.

\emph{Necessity.}
Suppose instead that $B\ge M+1$. Put $u=\max(s,t)$ and
\begin{equation}\label{eq:witness}
 n_*=rM+u.
\end{equation}
For the first count in \eqref{eq:counts}, the lower endpoint is
$rC+u-s+1$, strictly between $rC$ and $r(C+1)$, and the upper endpoint
is $rM+u<r(M+1)$. Consequently its indices are $C+1,\ldots,M$,
giving $A$ choices, with an empty range when $A=0$.
For the second count, the lower endpoint is
$u-s-t+1=1-\min(s,t)\le0$, and the upper endpoint is $rA+u-t$,
where $0\le u-t<r$. Its indices are $0,\ldots,A$, giving $A+1$ choices.
Neither count is clipped above, since $B\ge M+1$. Therefore
\begin{equation}\label{eq:drop}
 p_{n_*}-p_{n_*-1}=A-(A+1)=-1.
\end{equation}
This drop is in the first half: $n_*\ge1$ and
\[
 D-2n_*=r(B-M)-|s-t|-2\ge r-(r-2)-2=0.
\]
It contradicts \eqref{eq:half}. The alternatives $B\le M$ and $B\ge M+1$
exhaust the nondivisible case and complete the proof.
\end{proof}

Zero quotients $A$ or $C$ were allowed throughout. In particular, $a=1$
recovers the one-factor criterion, and $a=c=1$ gives unimodality exactly
when $b=1$. Thus no small-length exceptions are suppressed.

\section{Checks, an example, and scope}
For example, the theorem makes $[7]_q[9]_q[3]_{q^5}$ unimodal.
Its coefficients through the midpoint, at degrees $0$ through $12$, are
\[
 1,2,3,4,5,7,9,10,11,11,12,13,13.
\]
Both $[7]_q[3]_{q^5}$ and $[9]_q[3]_{q^5}$ fail the one-factor criterion
\cite[Proposition~5.3]{CIMSY2026}; applying that criterion to either
subproduct therefore does not account for the two-factor range.

Exact integer computations checked all $81\,648$ ordered tuples with
\[
 2\le r\le19,\qquad 1\le a,c\le18,\qquad 1\le b\le14,
\]
as well as $600$ draws with $2\le r\le79$, $1\le a,c\le120$, and
$1\le b\le20$. Python's \texttt{random.Random} was initialized once with seed
$202609140919$; each draw used four successive inclusive \texttt{randint}
calls in the order $(r,a,c,b)$, with repeated tuples permitted.
All $82\,248$ tests agreed with \eqref{eq:criterion}. The $54\,281$
nonunimodal tests also satisfied the exact witness identity
\eqref{eq:drop}.

For reproducibility, the coefficient construction first forms
\[
 t_\ell=\#\{(i,j):0\le i<a,\ 0\le j<c,\ i+j=\ell\},
 \qquad p_n=\sum_{j=0}^{b-1}t_{n-rj},
\]
with $t_\ell=0$ outside its support. The complete coefficient sequence is
scanned for a rise after a strict fall, without using the midpoint
criterion as the test. A separately implemented evaluator sums the three
affine pieces of the ordinary trapezoid over arithmetic progressions;
it agreed on $597\,941$ selected coefficient comparisons.
Reconstructing the objects and checking integer predicates with separate
implementations is also the verification design used by Zeng, Liu,
Ratnavelu, and Ong in their finite geometric certificates
\cite[Sections~4 and~6]{ZengEtAl2026}. Here the computations serve only as
consistency checks on the symbolic proof; their finite coverage is not
an argument for untested parameters.

The theorem gives an exact decision from the four input integers and a
degree witnessing failure. The all-$k$ sufficiency and the remaining
necessity assertions of the original conjecture remain unresolved here.

\paragraph{Computational and AI assistance.}
OpenAI Codex assisted with proof development, code, source comparison, and
drafting. A separate model context examined the proof and ran separately
written exact checks. This checking pass was not external peer review or
statistically independent evidence. No proof assistant was used.

\bibliographystyle{plain}
\small
\bibliography{references}
\normalsize
\end{document}
